\documentclass[10pt,letterpaper]{article}
\usepackage{cornell-notes}
\usepackage{mathtools}

\title{Chapter 3: Interpolation and Extrapolation}
\author{}
\date{}

\setDocTitle{Chapter 3: Interpolation and Extrapolation}
\setDocSubtitle{Numerical Methods}
\setDocVersion{v1.0}
\setDocStatus{Study Companion}
\setDocOwner{Numerical Methods Cornell Notes}
\setDocAuthor{}
\setDocDate{\today}
\setDocSummary{Cornell-format study and review notes for Chapter 3: Interpolation and Extrapolation.}

\setCornellCollection{Numerical Methods Cornell Notes}
\setCornellUnitType{Chapter}
\setCornellUnitNumber{3}
\setCornellUnitTitle{Interpolation and Extrapolation}
\setCornellCompanionLabel{Study and review companion}

\hypersetup{pdftitle={Chapter 3: Interpolation and Extrapolation},pdfsubject={Cornell notes},pdfauthor={}}

\begin{document}
\maketitle

\section*{Chapter Overview}
This chapter organizes the mathematical and computational ideas behind interpolation and extrapolation. The notes preserve the numbered section sequence while emphasizing method selection, explicit assumptions, numerical reliability, cost, and verification.

\section*{Learning Objectives}
Explain the governing problem class; compare Preliminaries: Searching an Ordered Table, Polynomial Interpolation and Extrapolation, Cubic Spline Interpolation, and Rational Function Interpolation and Extrapolation; design defensible stopping and verification checks; distinguish problem conditioning from algorithmic stability.

\section*{Prerequisites}
Polynomials, divided differences, linear systems, and multivariable coordinates.

\section*{Operational Workflow}
Locate a neighborhood; select a basis and degree consistent with smoothness and geometry; evaluate with a stable form; estimate local error; reject unsafe extrapolation.

\subsection*{Formula and notation anchor}
\[
p(x)=\sum_{i=0}^{n} y_i\,L_i(x),\qquad L_i(x)=\prod_{j\ne i}\frac{x-x_j}{x_i-x_j}
\]
The Lagrange form establishes the interpolant; barycentric or divided-difference forms are normally preferable for evaluation.

\begin{warnbox}{Numerical warning}
Runge oscillation, repeated or clustered nodes, extrapolation, spurious rational poles, poorly scaled coordinates, and dimensionality-driven cost.
\end{warnbox}

\section{3.0 Introduction}
\begin{cornell}
\noterow{What is the central idea?}{Interpolation estimates values between samples; extrapolation extends beyond them and is normally more sensitive to modeling error and noise.}
\noterow{How should it be used?}{For Introduction, begin from explicit inputs, outputs, preconditions, and representation choices.  Locate a neighborhood; select a basis and degree consistent with smoothness and geometry; evaluate with a stable form; estimate local error; reject unsafe extrapolation.}
\noterow{Cost and reliability}{Analyze arithmetic work, storage, data movement, and convergence separately.  Relevant failure pressures include runge oscillation, repeated or clustered nodes, extrapolation, spurious rational poles, poorly scaled coordinates, and dimensionality-driven cost.  Use scaled tolerances and report both success criteria and diagnostic evidence.}
\noterow{Implementation checklist}{\begin{itemize}
\item Validate dimensions, domains, ordering, and structural assumptions.
\item Guard small divisors, invalid states, overflow, underflow, and nonfinite values.
\item Make mutation, workspace, indexing, and termination behavior explicit.
\item Test a simple exact case, a difficult valid case, a boundary case, and an expected failure.
\end{itemize}}
\end{cornell}
\begin{summarybox}{Section summary}
Interpolation estimates values between samples; extrapolation extends beyond them and is normally more sensitive to modeling error and noise.  Select this technique only when its assumptions match the data, and accept a result only after an independent residual, error estimate, invariant, or refinement check.
\end{summarybox}

\section{3.1 Preliminaries: Searching an Ordered Table}
\begin{cornell}
\noterow{What is the central idea?}{Efficient interpolation begins by locating the correct interval while preserving bracketing invariants and handling boundaries and repeated abscissas.}
\noterow{How should it be used?}{For Preliminaries: Searching an Ordered Table, begin from explicit inputs, outputs, preconditions, and representation choices.  Locate a neighborhood; select a basis and degree consistent with smoothness and geometry; evaluate with a stable form; estimate local error; reject unsafe extrapolation.}
\noterow{Quantitative lens}{Identify the governing equation, recurrence, probability law, or geometric predicate for Preliminaries: Searching an Ordered Table.  Define every variable and scale; then derive a residual, error estimate, or invariant that can be checked independently.}
\noterow{Cost and reliability}{Analyze arithmetic work, storage, data movement, and convergence separately.  Relevant failure pressures include runge oscillation, repeated or clustered nodes, extrapolation, spurious rational poles, poorly scaled coordinates, and dimensionality-driven cost.  Use scaled tolerances and report both success criteria and diagnostic evidence.}
\noterow{Implementation checklist}{\begin{itemize}
\item Validate dimensions, domains, ordering, and structural assumptions.
\item Guard small divisors, invalid states, overflow, underflow, and nonfinite values.
\item Make mutation, workspace, indexing, and termination behavior explicit.
\item Test a simple exact case, a difficult valid case, a boundary case, and an expected failure.
\end{itemize}}
\end{cornell}
\begin{summarybox}{Section summary}
Efficient interpolation begins by locating the correct interval while preserving bracketing invariants and handling boundaries and repeated abscissas.  Select this technique only when its assumptions match the data, and accept a result only after an independent residual, error estimate, invariant, or refinement check.
\end{summarybox}

\section{3.2 Polynomial Interpolation and Extrapolation}
\begin{cornell}
\noterow{What is the central idea?}{A polynomial can match a finite set of data exactly, but degree, node placement, noise, and distance from the nodes govern usefulness.}
\noterow{How should it be used?}{For Polynomial Interpolation and Extrapolation, begin from explicit inputs, outputs, preconditions, and representation choices.  Locate a neighborhood; select a basis and degree consistent with smoothness and geometry; evaluate with a stable form; estimate local error; reject unsafe extrapolation.}
\noterow{Quantitative lens}{Identify the governing equation, recurrence, probability law, or geometric predicate for Polynomial Interpolation and Extrapolation.  Define every variable and scale; then derive a residual, error estimate, or invariant that can be checked independently.}
\noterow{Cost and reliability}{Analyze arithmetic work, storage, data movement, and convergence separately.  Relevant failure pressures include runge oscillation, repeated or clustered nodes, extrapolation, spurious rational poles, poorly scaled coordinates, and dimensionality-driven cost.  Use scaled tolerances and report both success criteria and diagnostic evidence.}
\noterow{Implementation checklist}{\begin{itemize}
\item Validate dimensions, domains, ordering, and structural assumptions.
\item Guard small divisors, invalid states, overflow, underflow, and nonfinite values.
\item Make mutation, workspace, indexing, and termination behavior explicit.
\item Test a simple exact case, a difficult valid case, a boundary case, and an expected failure.
\end{itemize}}
\end{cornell}
\begin{summarybox}{Section summary}
A polynomial can match a finite set of data exactly, but degree, node placement, noise, and distance from the nodes govern usefulness.  Select this technique only when its assumptions match the data, and accept a result only after an independent residual, error estimate, invariant, or refinement check.
\end{summarybox}

\section{3.3 Cubic Spline Interpolation}
\begin{cornell}
\noterow{What is the central idea?}{A cubic spline joins low-degree pieces with continuous derivatives, trading global oscillation for a tridiagonal solve and explicit boundary conditions.}
\noterow{How should it be used?}{For Cubic Spline Interpolation, begin from explicit inputs, outputs, preconditions, and representation choices.  Locate a neighborhood; select a basis and degree consistent with smoothness and geometry; evaluate with a stable form; estimate local error; reject unsafe extrapolation.}
\noterow{Quantitative anchor}{\[
s_i(x_i)=y_i,\quad s_i'(x_{i+1})=s_{i+1}'(x_{i+1})
\]
State the dimensions, scaling, and validity assumptions before using this relation.  Check the resulting residual or invariant rather than trusting an iteration count alone.}
\noterow{Cost and reliability}{Analyze arithmetic work, storage, data movement, and convergence separately.  Relevant failure pressures include runge oscillation, repeated or clustered nodes, extrapolation, spurious rational poles, poorly scaled coordinates, and dimensionality-driven cost.  Use scaled tolerances and report both success criteria and diagnostic evidence.}
\noterow{Implementation checklist}{\begin{itemize}
\item Validate dimensions, domains, ordering, and structural assumptions.
\item Guard small divisors, invalid states, overflow, underflow, and nonfinite values.
\item Make mutation, workspace, indexing, and termination behavior explicit.
\item Test a simple exact case, a difficult valid case, a boundary case, and an expected failure.
\end{itemize}}
\end{cornell}
\begin{summarybox}{Section summary}
A cubic spline joins low-degree pieces with continuous derivatives, trading global oscillation for a tridiagonal solve and explicit boundary conditions.  Select this technique only when its assumptions match the data, and accept a result only after an independent residual, error estimate, invariant, or refinement check.
\end{summarybox}

\section{3.4 Rational Function Interpolation and Extrapolation}
\begin{cornell}
\noterow{What is the central idea?}{Ratios of polynomials can represent poles and steep behavior more naturally than a single polynomial, but spurious denominator zeros require detection.}
\noterow{How should it be used?}{For Rational Function Interpolation and Extrapolation, begin from explicit inputs, outputs, preconditions, and representation choices.  Locate a neighborhood; select a basis and degree consistent with smoothness and geometry; evaluate with a stable form; estimate local error; reject unsafe extrapolation.}
\noterow{Quantitative lens}{Identify the governing equation, recurrence, probability law, or geometric predicate for Rational Function Interpolation and Extrapolation.  Define every variable and scale; then derive a residual, error estimate, or invariant that can be checked independently.}
\noterow{Cost and reliability}{Analyze arithmetic work, storage, data movement, and convergence separately.  Relevant failure pressures include runge oscillation, repeated or clustered nodes, extrapolation, spurious rational poles, poorly scaled coordinates, and dimensionality-driven cost.  Use scaled tolerances and report both success criteria and diagnostic evidence.}
\noterow{Implementation checklist}{\begin{itemize}
\item Validate dimensions, domains, ordering, and structural assumptions.
\item Guard small divisors, invalid states, overflow, underflow, and nonfinite values.
\item Make mutation, workspace, indexing, and termination behavior explicit.
\item Test a simple exact case, a difficult valid case, a boundary case, and an expected failure.
\end{itemize}}
\end{cornell}
\begin{summarybox}{Section summary}
Ratios of polynomials can represent poles and steep behavior more naturally than a single polynomial, but spurious denominator zeros require detection.  Select this technique only when its assumptions match the data, and accept a result only after an independent residual, error estimate, invariant, or refinement check.
\end{summarybox}

\section{3.5 Coefficients of the Interpolating Polynomial}
\begin{cornell}
\noterow{What is the central idea?}{Coefficient recovery converts value data into a chosen polynomial basis; the basis and node scaling strongly affect numerical conditioning.}
\noterow{How should it be used?}{For Coefficients of the Interpolating Polynomial, begin from explicit inputs, outputs, preconditions, and representation choices.  Locate a neighborhood; select a basis and degree consistent with smoothness and geometry; evaluate with a stable form; estimate local error; reject unsafe extrapolation.}
\noterow{Quantitative lens}{Identify the governing equation, recurrence, probability law, or geometric predicate for Coefficients of the Interpolating Polynomial.  Define every variable and scale; then derive a residual, error estimate, or invariant that can be checked independently.}
\noterow{Cost and reliability}{Analyze arithmetic work, storage, data movement, and convergence separately.  Relevant failure pressures include runge oscillation, repeated or clustered nodes, extrapolation, spurious rational poles, poorly scaled coordinates, and dimensionality-driven cost.  Use scaled tolerances and report both success criteria and diagnostic evidence.}
\noterow{Implementation checklist}{\begin{itemize}
\item Validate dimensions, domains, ordering, and structural assumptions.
\item Guard small divisors, invalid states, overflow, underflow, and nonfinite values.
\item Make mutation, workspace, indexing, and termination behavior explicit.
\item Test a simple exact case, a difficult valid case, a boundary case, and an expected failure.
\end{itemize}}
\end{cornell}
\begin{summarybox}{Section summary}
Coefficient recovery converts value data into a chosen polynomial basis; the basis and node scaling strongly affect numerical conditioning.  Select this technique only when its assumptions match the data, and accept a result only after an independent residual, error estimate, invariant, or refinement check.
\end{summarybox}

\section{3.6 Interpolation on a Grid in Multidimensions}
\begin{cornell}
\noterow{What is the central idea?}{Tensor-product interpolation reuses one-dimensional procedures along coordinate directions, with cost and storage growing rapidly with dimension.}
\noterow{How should it be used?}{For Interpolation on a Grid in Multidimensions, begin from explicit inputs, outputs, preconditions, and representation choices.  Locate a neighborhood; select a basis and degree consistent with smoothness and geometry; evaluate with a stable form; estimate local error; reject unsafe extrapolation.}
\noterow{Quantitative lens}{Identify the governing equation, recurrence, probability law, or geometric predicate for Interpolation on a Grid in Multidimensions.  Define every variable and scale; then derive a residual, error estimate, or invariant that can be checked independently.}
\noterow{Cost and reliability}{Analyze arithmetic work, storage, data movement, and convergence separately.  Relevant failure pressures include runge oscillation, repeated or clustered nodes, extrapolation, spurious rational poles, poorly scaled coordinates, and dimensionality-driven cost.  Use scaled tolerances and report both success criteria and diagnostic evidence.}
\noterow{Implementation checklist}{\begin{itemize}
\item Validate dimensions, domains, ordering, and structural assumptions.
\item Guard small divisors, invalid states, overflow, underflow, and nonfinite values.
\item Make mutation, workspace, indexing, and termination behavior explicit.
\item Test a simple exact case, a difficult valid case, a boundary case, and an expected failure.
\end{itemize}}
\end{cornell}
\begin{summarybox}{Section summary}
Tensor-product interpolation reuses one-dimensional procedures along coordinate directions, with cost and storage growing rapidly with dimension.  Select this technique only when its assumptions match the data, and accept a result only after an independent residual, error estimate, invariant, or refinement check.
\end{summarybox}

\section{3.7 Interpolation on Scattered Data in Multidimensions}
\begin{cornell}
\noterow{What is the central idea?}{Scattered-data interpolation requires a neighborhood or triangulation model because no regular cell structure identifies the contributing samples.}
\noterow{How should it be used?}{For Interpolation on Scattered Data in Multidimensions, begin from explicit inputs, outputs, preconditions, and representation choices.  Locate a neighborhood; select a basis and degree consistent with smoothness and geometry; evaluate with a stable form; estimate local error; reject unsafe extrapolation.}
\noterow{Quantitative lens}{Identify the governing equation, recurrence, probability law, or geometric predicate for Interpolation on Scattered Data in Multidimensions.  Define every variable and scale; then derive a residual, error estimate, or invariant that can be checked independently.}
\noterow{Cost and reliability}{Analyze arithmetic work, storage, data movement, and convergence separately.  Relevant failure pressures include runge oscillation, repeated or clustered nodes, extrapolation, spurious rational poles, poorly scaled coordinates, and dimensionality-driven cost.  Use scaled tolerances and report both success criteria and diagnostic evidence.}
\noterow{Implementation checklist}{\begin{itemize}
\item Validate dimensions, domains, ordering, and structural assumptions.
\item Guard small divisors, invalid states, overflow, underflow, and nonfinite values.
\item Make mutation, workspace, indexing, and termination behavior explicit.
\item Test a simple exact case, a difficult valid case, a boundary case, and an expected failure.
\end{itemize}}
\end{cornell}
\begin{summarybox}{Section summary}
Scattered-data interpolation requires a neighborhood or triangulation model because no regular cell structure identifies the contributing samples.  Select this technique only when its assumptions match the data, and accept a result only after an independent residual, error estimate, invariant, or refinement check.
\end{summarybox}

\section{3.8 Laplace Interpolation}
\begin{cornell}
\noterow{What is the central idea?}{Laplace interpolation fills missing values by enforcing a discrete harmonic condition, turning local averaging into a sparse boundary-value problem.}
\noterow{How should it be used?}{For Laplace Interpolation, begin from explicit inputs, outputs, preconditions, and representation choices.  Locate a neighborhood; select a basis and degree consistent with smoothness and geometry; evaluate with a stable form; estimate local error; reject unsafe extrapolation.}
\noterow{Quantitative lens}{Identify the governing equation, recurrence, probability law, or geometric predicate for Laplace Interpolation.  Define every variable and scale; then derive a residual, error estimate, or invariant that can be checked independently.}
\noterow{Cost and reliability}{Analyze arithmetic work, storage, data movement, and convergence separately.  Relevant failure pressures include runge oscillation, repeated or clustered nodes, extrapolation, spurious rational poles, poorly scaled coordinates, and dimensionality-driven cost.  Use scaled tolerances and report both success criteria and diagnostic evidence.}
\noterow{Implementation checklist}{\begin{itemize}
\item Validate dimensions, domains, ordering, and structural assumptions.
\item Guard small divisors, invalid states, overflow, underflow, and nonfinite values.
\item Make mutation, workspace, indexing, and termination behavior explicit.
\item Test a simple exact case, a difficult valid case, a boundary case, and an expected failure.
\end{itemize}}
\end{cornell}
\begin{summarybox}{Section summary}
Laplace interpolation fills missing values by enforcing a discrete harmonic condition, turning local averaging into a sparse boundary-value problem.  Select this technique only when its assumptions match the data, and accept a result only after an independent residual, error estimate, invariant, or refinement check.
\end{summarybox}

\section*{Method comparison}
\begin{center}
\small
\begin{tabularx}{\linewidth}{>{\bfseries\raggedright\arraybackslash}p{0.22\linewidth}XX}
\toprule
Method or lens & Best use & Main caution \\
\midrule
Local polynomial & Smooth local data & Degree and node placement control oscillation \\
Cubic spline & Smooth tabulated curves & Requires boundary conditions \\
Rational & Pole-like or steep behavior & Denominator can create spurious poles \\
\bottomrule
\end{tabularx}
\end{center}

\section*{Worked example}
\begin{exambox}{Independent worked example}
For samples $(0,1)$, $(1,2)$, and $(2,5)$, the quadratic interpolant is $p(x)=x^2+1$, so $p(1.5)=3.25$.  Substitution at all three nodes gives zero interpolation residual.  This exact fit does not estimate noise; clustered nodes or extrapolation beyond $[0,2]$ can amplify error.
\end{exambox}

\section*{Chapter synthesis}
\begin{summarybox}{Chapter synthesis}
The chapter's methods should be viewed as a decision system rather than a list of routines.  Begin with the mathematical problem and its conditioning, exploit structure, select an algorithm with compatible stability and cost, and verify the computed answer with evidence that is independent of the internal iteration.  The recurring implementation discipline is: locate a neighborhood; select a basis and degree consistent with smoothness and geometry; evaluate with a stable form; estimate local error; reject unsafe extrapolation.
\end{summarybox}

\subsection*{Common mistakes and numerical hazards}
\begin{warnbox}{Common mistakes and numerical hazards}
Runge oscillation, repeated or clustered nodes, extrapolation, spurious rational poles, poorly scaled coordinates, and dimensionality-driven cost.  A second recurring mistake is reporting only a computed value: always retain tolerances, iteration counts, residuals, refinement behavior, and relevant structural checks.
\end{warnbox}

\subsection*{Retrieval practice}
\textbf{Questions}
\begin{enumerate}
\item What problem does Preliminaries: Searching an Ordered Table address, and which assumption is easiest to violate?
\item What problem does Polynomial Interpolation and Extrapolation address, and which assumption is easiest to violate?
\item What problem does Cubic Spline Interpolation address, and which assumption is easiest to violate?
\item What problem does Interpolation on Scattered Data in Multidimensions address, and which assumption is easiest to violate?
\item What problem does Laplace Interpolation address, and which assumption is easiest to violate?
\item How do conditioning, stability, and the final residual play different roles in this chapter?
\end{enumerate}

\textbf{Brief answers}
\begin{enumerate}
\item Efficient interpolation begins by locating the correct interval while preserving bracketing invariants and handling boundaries and repeated abscissas. Recheck the section's domain, structure, scaling, and failure diagnostics.
\item A polynomial can match a finite set of data exactly, but degree, node placement, noise, and distance from the nodes govern usefulness. Recheck the section's domain, structure, scaling, and failure diagnostics.
\item A cubic spline joins low-degree pieces with continuous derivatives, trading global oscillation for a tridiagonal solve and explicit boundary conditions. Recheck the section's domain, structure, scaling, and failure diagnostics.
\item Scattered-data interpolation requires a neighborhood or triangulation model because no regular cell structure identifies the contributing samples. Recheck the section's domain, structure, scaling, and failure diagnostics.
\item Laplace interpolation fills missing values by enforcing a discrete harmonic condition, turning local averaging into a sparse boundary-value problem. Recheck the section's domain, structure, scaling, and failure diagnostics.
\item Conditioning describes sensitivity of the mathematical problem, stability describes error propagation by the algorithm, and the residual measures satisfaction of the computed equations; none is a universal substitute for the others.
\end{enumerate}

\subsection*{Mastery checklist}
\begin{itemize}
\item[$\square$] I can explain and select Preliminaries: Searching an Ordered Table without relying on a memorized code listing.
\item[$\square$] I can explain and select Polynomial Interpolation and Extrapolation without relying on a memorized code listing.
\item[$\square$] I can explain and select Cubic Spline Interpolation without relying on a memorized code listing.
\item[$\square$] I can explain and select Rational Function Interpolation and Extrapolation without relying on a memorized code listing.
\item[$\square$] I can explain and select Coefficients of the Interpolating Polynomial without relying on a memorized code listing.
\item[$\square$] I can state a scaled stopping rule and an independent verification check.
\item[$\square$] I can identify at least one conditioning risk and one algorithmic stability risk.
\item[$\square$] I can estimate time, storage, and data-movement costs at the appropriate level.
\end{itemize}

\end{document}

